% TC33 — SOS มีผู้รับฝั่งสตาฟ
\begin{testscenario}{33}{4}
  \tcmeta{TC33}{CS 2568/01}{ระบบเว็บแอปพลิเคชัน ``NUCARE''}{นายสมชาย ใจดี}{วันที่ 7 พฤศจิกายน พ.ศ.\ 2569}{1.0}
  \tcname{การแจ้งช่วยเหลือด่วน (SOS) มีผู้รับและดำเนินการได้}
  \tcpre{บัญชีนิสิต มีที่ปรึกษา · บัญชี MENTAL\_HEALTH · รัน unit emergency/SOS}
  \tcdesc{ตรวจว่าเมื่อนิสิตกด SOS ระบบแจ้งทีมสุขภาพจิตและที่ปรึกษา และผู้เชี่ยวชาญรับ/ปิดเคสได้}

  \tccase{1}{ลำดับสถานะ SOS}{%
    \begin{tcsteps}
      \item รัน \tccode{emergency-case-status.test.ts}
    \end{tcsteps}}{%
    OPEN$\to$IN\_PROGRESS/RESOLVED ได้ · ปิดแล้วเปลี่ยนต่อไม่ได้}{%
    ECS-01..04 ผ่าน}{Pass}

  \tccase{2}{รวมผู้รับแจ้งเตือน}{%
    \begin{tcsteps}
      \item รัน \tccode{notify-sos-recipients.test.ts}
    \end{tcsteps}}{%
    รวม MH + ที่ปรึกษา และตัดซ้ำ}{%
    SOS-R01 ผ่าน}{Pass}

  \tccase{3}{นิสิตกด SOS แล้วทีมเห็นคิว}{%
    \begin{tcsteps}
      \item เข้าสู่ระบบนิสิตแล้วกด SOS
      \item เข้า \tccode{mh.test} ดูหน้าหลัก
    \end{tcsteps}}{%
    มีแจ้งเตือน SOS\_OPENED และรายการคำขอช่วยเหลือด่วน}{%
    API สร้าง EmergencyCase + alerts และแดชบอร์ดโหลด emergencyCases}{Pass}

  \tccase{4}{รับดำเนินการแล้วปิดเคส}{%
    \begin{tcsteps}
      \item กดรับดำเนินการแล้วปิดเคสจากหน้าหลัก MH
    \end{tcsteps}}{%
    สถานะ IN\_PROGRESS แล้ว RESOLVED · แจ้งนิสิต · ปิด alerts เมื่อปิดเคส}{%
    PATCH emergency-cases/:id/status ทำงานตามลำดับ}{Pass}
\end{testscenario}
